\scanpage{293}
\mspage{17}
% archive sheet number: 292

% slanted
\margin{Je vais désigner \deleted{\unsure{Loc|par}} par Loc 6 l'axiome Loc 6(a,b), sans inclure c.}

\uline{NB} \unsure{1} Je distingue ces conditions comme $\mathrm{Loc}(6a)$, \deleted{\unsure{et}} $\mathrm{Loc}(6b)$, \inserted{$\mathrm{Loc}(6c)$} Je présume qu'on n'aura jamais à faire usage \unsure{à|de} $\mathrm{Loc}(6b)$, sans supposer également $\mathrm{Loc}(6a)$, mais l'inverse n'est sûrement pas le cas. Ainsi $W_0=\{f\in\unsure{\mathrm{Fl}|\mathrm{Fl.}}\mathrm{Cat}\mid\pi_0(f)\ \text{bijectif}\}$ satisfait a), mais non b). Je présume (sans avoir vérifié) que c'est pareil pour les $W_n$, $n\in\unsure{N|\mathbb N}$, de la ``\uline{$n$-équivalence homotopique}''.

Il se pourrait que $\uline{Loc}(6a)$ et $\uline{Loc}(6a,b)$ soient \unsure{équivalentes|équivalents} à la forme plus faible, et plus simple, suivante. (Si oui, c'est sous la forme la plus simple, bien sûr, qu'il convient de les énoncer.)

\uline{Loc} \deleted{\unsure{6}}$6'$ (Axiome des carrés cocartésiens, \inserted{ou de Mayer-Vietoris}).

Soit
\diagram{\begin{tikzcd}[ampersand replacement=\&,column sep=small]
X_0 \arrow[rr,"i_2"] \arrow[dr,"i_1"'] \& \& X_2\arrow[dr,"j_1"] \&\\
\& X_1\arrow[rr,"j_2"'] \& \& X_3
\end{tikzcd}}
un diagramme \deleted{\unsure{de}} $\mathrm{MV}$ dans $\mathrm{Cat}$. Alors
\[
\begin{aligned}
\text{a)}\quad&\text{Si } i_1\in W,\quad j_1\in W.\\
\text{b)}\quad&\text{Si } j_1\in W,\quad i_1\in W.
\end{aligned}
\]
